% =============================================================================
%  config/theorems.tex: ambienti per definizioni, teoremi, esempi, ...
%
%  Uso:  \begin{teorema}[Nome opzionale]\label{thm:...} ... \end{teorema}
%        \begin{proof} ... \end{proof}     (titolo "Dimostrazione" via babel)
%
%  Numerazione condivisa per capitolo: Definizione 2.1, Teorema 2.2, ...
% =============================================================================

% --- Stili ----------------------------------------------------------------------
\declaretheoremstyle[
  headfont = \bfseries\color{accento},
  bodyfont = \itshape,
  spaceabove = 0.8em, spacebelow = 0.8em,
]{enunciato}

\declaretheoremstyle[
  headfont = \bfseries\color{accento},
  bodyfont = \normalfont,
  spaceabove = 0.8em, spacebelow = 0.8em,
]{definizione}

\declaretheoremstyle[
  headfont = \bfseries\itshape,
  bodyfont = \normalfont,
  spaceabove = 0.6em, spacebelow = 0.6em,
]{nota}

% --- Ambienti -------------------------------------------------------------------
% Il contatore "teorema" è condiviso da tutti gli ambienti numerati.
\declaretheorem[style=enunciato,   name=Teorema,      numberwithin=chapter]{teorema}
\declaretheorem[style=enunciato,   name=Lemma,        sibling=teorema]{lemma}
\declaretheorem[style=enunciato,   name=Proposizione, sibling=teorema]{proposizione}
\declaretheorem[style=enunciato,   name=Corollario,   sibling=teorema]{corollario}
\declaretheorem[style=definizione, name=Definizione,  sibling=teorema]{definizione}
\declaretheorem[style=definizione, name=Esempio,      sibling=teorema]{esempio}
\declaretheorem[style=definizione, name=Esercizio,    sibling=teorema]{esercizio}
\declaretheorem[style=nota,        name=Osservazione, sibling=teorema]{osservazione}
\declaretheorem[style=nota,        name=Soluzione,    numbered=no]{soluzione}

% --- Nomi per \cref (singolare/plurale) ---------------------------------------------
\crefname{teorema}{teorema}{teoremi}
\crefname{lemma}{lemma}{lemmi}
\crefname{proposizione}{proposizione}{proposizioni}
\crefname{corollario}{corollario}{corollari}
\crefname{definizione}{definizione}{definizioni}
\crefname{esempio}{esempio}{esempi}
\crefname{esercizio}{esercizio}{esercizi}
\crefname{osservazione}{osservazione}{osservazioni}

% --- Algoritmi ------------------------------------------------------------------------
% Riquadro per la descrizione a passi degli algoritmi (simplesso, B&B, ...).
\newtcolorbox{algoritmo}[1]{
  enhanced, breakable,
  colback = accento!4, colframe = accento,
  fonttitle = \bfseries, title = {Algoritmo: #1},
}
